%%% ALLGEMEINE FORMATIERUNGEN
\documentclass[%
	twoside=false,% Doppelseitig schreiben oder nicht. Wichtig für automatische Bindekorrektur (BCOR).
	paper=a4,% Papiergröße festlegen, hier A4
	DIV=15,% Mit Hilfe der Option DIV=Faktor wird festgelegt, in wie viele Streifen die Seite horizontal und vertikal bei der Satzspiegelkonstruktion eingeteilt wird. DIV=calc nimmt einem diese Arbeit ab (hängt ja auch von der Schriftart und -größe ab usw.)
	BCOR=0mm,% Mit Hilfe der Option BCOR=Korrektur geben Sie den absoluten Wert der Bindekorrektur an, also die Breite des Bereichs, der durch die Bindung von der Papierbreite verloren geht. Am besten in der Druckerei nachfragen! Man kann einen Wert von ca. 10mm für Arbeiten bis 100 Seiten annehmen.
	parskip=false,% Wie Absätze gekennzeichnet werden. Zitat: "parskip=false Absätze werden durch einen Einzug der ersten Zeilen von einem Geviert (1 em) gekennzeichnet. Der erste Absatz eines Abschnitts wird nicht eingezogen." Alles andere ist hier nicht sinnvoll, siehe scrguide.pdf und Leitfaden_fuer_wiss_Dokumente.pdf!
	fontsize=11pt,% Schriftgröße
	listof=totoc,% Abbildungsverz. usw. werden in das Inhaltsverzeichnis aufgenommen
	bibliography=totoc,%  Literaturverzeichnis wird in das Inhaltsverzeichnis aufgenommen
	captions=abovetabular% Die Beschriftungen für Tabellen sind über der Tabelle
	]{scrreprt}%
%
%
%%% HEADER und FOOTER DESIGN
\usepackage[% Header und Footer designen
	autooneside=true,% ???
	headsepline=true,% Trennungsstrich für den Header (außer bei neuen Kapiteln und bei den Indizes)
	footsepline=true,% Trennungsstrich für den Footer (außer bei neuen Kapiteln und bei den Indizes)
	plainfootsepline=true% Trennungsstrich für den Footer bei neuen Kapiteln und bei den Indizes
	]{scrlayer-scrpage}%
%
%
%%% WAS steht WO im HEADER und FOOTER
\pagestyle{scrheadings}% scrheadings einstellen, damit man alles weitere im Header und Footer richtig definieren kann
\clearscrheadfoot% Header und footer clearen (sonst hat man die Seitenzahl eventuell zwei mal im Footer...)
\automark[section]{chapter}% Von den ersten beiden Gliederungsebenen jeweils die niedrigste im Header anzeigen
\chead{\headmark}% Zentrierte Überschrift im Header
 \ofoot[\pagemark]{\pagemark}% Seitenzahl für einseitigen Druck im Footer auf der rechten Seite platzieren, für zweiseitigen Druck immer außen platzieren
\addtokomafont{pagefoot}{\small}% Fußnoten etwas kleiner als die normale Schriftgröße setzen - zur Sicherheit den Befehl wiederholen - falls sich die Standard-Einstellungen ändern (safety first!)
%
%
%%% SPRACHE
\usepackage{polyglossia}% statt Babel nutzt man für xelatex polyglossia!
\setmainlanguage[%
	variant=german,%es geht auch austrian oder swiss
	spelling=new,%neue Rechtschreibung ein/aus
	latesthyphen=true,%ausschalten bei bugs - neueste hyphenation
	babelshorthands=false%Trennung wie beim Paket Babel ausschalten		
	]%
{german}%
%
%
\usepackage[T1]{fontenc}% Damit man Text sinnvoll kopieren kann (u.a. Umlaute usw. - damit ein Ü richtig im Inhaltsverzeichnis der PDF angezeigt wird
%
%
%%% SCHRIFTART
\usepackage{fontspec}% "Normale" Schriftarten können verwendet werden (die, die auf dem PC selbst installiert sind)
\setmainfont{Arial}% Schriftart festlegen
%
%
%%% SPRACHE mit babel
% \usepackage[ngerman]{babel}% Für richtige Silben- und Worttrennung in der jeweiligen Sprache und die richtige Sprache in den automatisch generierten Verzeichnissen wird das Paket babel genutzt.
% Paket babel ngerman für die "neue" deutsche Sprache nutzen - eingeführt 1998, es ist 2018!
% siehe: https://tex.stackexchange.com/questions/67549/whats-the-difference-between-ngerman-and-german-in-babel
%
%
%%% ZEILENABSTAND
\usepackage{setspace}% Zeilenabstand einstellen mit setspace.sty - so werden auch in Fußzeilen die alten Abstände beibehalten! Siehe Sündenkatalog
\onehalfspacing% 1,5 Zeilig gewählt
% oder, falls dieses Paket mal nicht funktioniert:
% \linespread{1.2} % für Werte schaut man hier: https://www.sharelatex.com/learn/Paragraph_formatting#Line_spacing
%
%
%%% LEERE SEITEN
%\usepackage{afterpage}% Paket für eine leere Seite laden, damit unser Dokument genau so aussieht, wie gefordert (insbesondere nötig für zweiseitigen Druck!!!
% siehe https://tex.stackexchange.com/questions/36880/insert-a-blank-page-after-current-page
%\newcommand\blankpage{% Befehl neu definieren - leere Seite ohne Header und Footer und den Counter für das Seitenzählen zurücksetzen
	%\null%
    %\thispagestyle{empty}%
    %\addtocounter{page}{-1}%
    %\newpage}% %Alles was mit dem Format und Aufbau zu tun hat, ist hier zu finden
%
% 
\usepackage{amsmath} % Das einzig wahre Paket für tolle Formeln
%
%
\usepackage{graphicx} % Paket für sinnvolles einfügen von Grafiken
\graphicspath{ {./gfx/} } % Pfad für die Bilder festlegen - relativ zum Haupttex-file
%
%
% Das Paket listings wird benötigt, um line-breaks zu haben, wenn man normalen .tex-code eingibt
% weiteres dazu: https://tex.stackexchange.com/questions/121601/automatically-wrap-the-text-in-verbatim?noredirect=1&lq=1
\usepackage{listings}
\lstset{
basicstyle=\small\ttfamily,
columns=flexible,
breaklines=true
}
%
%
% Das Paket glossaries wird benötigt, um das Abkürzungsverzeichnis und das Verzeichnis für Formelzeichen und Symbole zu erstellen. Die Option toc heißt, dass diese im Inhaltsverzeichnis gelistet werden.
%\usepackage[toc,acronym]{glossaries}
% Mehr Optionen und Stile laden
\usepackage{glossaries-extra}
\usepackage{glossary-longbooktabs} %benötigt?
\usepackage{glossaries-extra-stylemods} %benötigt?
% Nun noch die Verzeichnisse erstellen - ohne index - muss verbessert werden
\makenoidxglossaries 
% neue Einträge erstellen
\newglossaryentry{formula}
{
    name=formula,
    description={A mathematical expression}
}
\newglossaryentry{latex}
{
    name=latex,
    description={Is a mark up language specially suited for 
scientific documents}
}
%
%
% tabu Paket, damit man eine Tabelle einfach so auf die Breite des Textfeldes einstellen kann.
% Alternativ tabularx?
\usepackage{tabu}
%
%
\usepackage{booktabs} % booktabs für schönere Tabellen
% benötigt?
%
%
\usepackage[bookmarks=true,bookmarksnumbered=true,hypertexnames=false,breaklinks=true,linkbordercolor={0 0 1},pdfborder={0 0 1}]{hyperref} % Mit dem Paket hyperref können Links innerhalb des Dokumentes im späteren PDF Format angeklickt werden und schicken einen durch das PDF, so wie es sein soll. Tolles Feature, sollten alle PDFs haben. Zusätzlich können für die verschiedenen Reader noch .. definiert werden%
% siehe http://www.pa.op.dlr.de/~PatrickJoeckel/pdflatex/index.html für eine tolle Hilfe
% Angaben in die PDF-Infos uebernehmen für PDF Reader
\makeatletter%
\hypersetup{ % PDF-Dokument-Informationen setzen
            pdftitle={\@title}, %Titel der Arbeit
            pdfauthor={\@author}, %Author
            %pdfkeywords={}, % CR-Klassifikation und ggf. weitere Stichworte
            %pdfsubject={} %
}%
\makeatother%
%
% "Satzspiegelberechnung mit demselben DIV -Argument, das beim letzten Aufruf an gegeben wurde, erneut durchführen." laut Anleitung des KOMA-Scriptes, also tun wir das auch:
%\KOMAoptions{DIV=last}
% benötigt?
%
%
%Dokument beginnen
\begin{document}
%folgender Grundsätzlicher Aufbau der Diplomarbeit wird gefordert - die mit /// markierten Zeilen sind Optional und wurden für das Template ausgelassen!
% Leerseite
% Aufgabenblatt der Hochschule (1x original, 1x Kopie) - wird hier nicht eingefügt, da diese Kopie in der Druckerei eingefügt werden muss!
% ///Spezifisches Titelblatt
% ///Sperrvermerk
% ///Kurzreferat/Vorwort/Abstract
% Inhaltsverzeichnis
% Abkürzungsverzeichnis
% Verzeichnis der Formelzeichen und Symbole
% Abbildungsverzeichnis
% Tabellenverzeichnis
% Inhalt des Dokuments
% Literatur- und Quellenverzeichnis
% ///Danksagung
% Eidesstattliche Erklärung
% ///Erklärung über Arbeitsanteile bei Gruppendiplomarbeiten
% Anlagenverzeichnis
% Anlagen
% Leerseite
%
%
%
%
%
\titlehead{\centering\includegraphics[width=7cm]{htw-logo} } %Bild des HTW-Logos
\subject{Diplomarbeit}
\title{Diplomarbeitsthema hier einfügen} %
\subtitle{Untertitel } %Untertitel der Diplomarbeit
\author{Hier könnte dein Name stehen} %eigenen Namen einfügen, dieser wird automatisch in die Eidesstattliche Erklärung übernommen
\date{\today } %heutiges Datum
\publishers{Verlag }
%\thanks{Fußnote} %hier ausgeblendet, da nicht sinnvoll

% 1. Leerseite
%\afterpage{\blankpage}
% ///Spezifisches Titelblatt
\maketitle
\newpage
% ///Sperrvermerk
% ///Kurzreferat/Vorwort/Abstract
%Inhaltsverzeichnis
% Seitenzahlen mit großen römischen Buchstaben beginnen
\pagenumbering{Roman}
\tableofcontents
\newpage
% Abkürzungsverzeichnis
\setglossarystyle{long3col}
\printnoidxglossaries
\newpage
% Verzeichnis der Formelzeichen und Symbole
Verzeichnis der Formelzeichen und Symbole
\newpage
% Abbildungsverzeichnis
\listoffigures
\newpage
% Tabellenverzeichnis
\listoftables
\newpage
% Inhalt des Dokuments
% Damit folgt immer die richtige Seitenzahl, auch bei zweiseitigem Druck! Siehe KOMA-Script!
\cleardoubleoddpage
% hier beginnt die Zählung der Seiten in arabischen Zahlen
\pagenumbering{arabic}
%
%
%
\chapter{Quick and dirty - Vorwort}%
\label{chap: vorwort}%
Ein paar kleine Dinge müsst ihr einstellen, \emph{bevor} ihr anfangt:%
\begin{itemize}%
\item Wollt ihr ein einseitige oder zweiseitige Abschlussarbeit? Wenn ihr euch für den beidseitigen Druck entschieden habt, dann stellt ihr in der Datei \textsc{/style/form-und-aufbau.tex} in \textsc{Zeile 3} \emph{twoside=true} und in \textsc{Zeile 16} \emph{autooneside=false}.%
\item Wie groß ist eure Bindekorrektur? Fragt in der Druckerei nach und stellt in \textsc{Zeile 6} in demselben Dokument den Wert unter \emph{BCOR} entsprechend ein.%
\end{itemize}%
Anschließend tragt ihr in der Datei \textsc{/content/titlepage.tex} in \textsc{Zeile 1-8} euren Namen usw. ein. Anschließend könnt ihr mit eurer Abschlussarbeit starten.%

Extrem Ungeduldige und erfahrenere Latex-Benutzer finden in der pdf \textsc{quickanddirty-cheatsheet} die wichtigsten Befehle und legen direkt los. Alle anderen lesen jetzt weiter.\\%%
%
%
\chapter{Struktur und Texteigenschaften}%
\section{Allgemeiner Aufbau}%
Folgender Grundsätzlicher Aufbau der Diplomarbeit wird gefordert - von den optionalen Seiten laut Leitfaden optional ist nur das "Spezifische Titelblatt" übernommen worden:%
\begin{itemize}%
 \item Leerseite%
 \item Aufgabenblatt der Hochschule (1x original, 1x Kopie) \emph{wird hier nicht eingefügt, da diese Kopie besser in der Druckerei eingefügt werden sollte!}%
 \item Spezifisches Titelblatt%
 \item Inhaltsverzeichnis%
 \item Abkürzungsverzeichnis%
 \item Verzeichnis der Formelzeichen und Symbole%
 \item Abbildungsverzeichnis%
 \item Tabellenverzeichnis%
 \item Inhalt des Dokuments%
 \item Literatur- und Quellenverzeichnis%
 \item Eidesstattliche Erklärung%
 \item Anlagenverzeichnis%
 \item Anlagen%
 \item Leerseite%
\end{itemize}%
%
%
\section{Allgemeine Texteingabe}%
Latex setzt den Text besser als die meisten Programme. \href{https://tex.stackexchange.com/questions/110133/visual-comparison-between-latex-and-word-output-hyphenation-typesetting-ligat#110140}{Hier} und \href{http://www.rtznet.nl/zink/latex.php}{hier} sind Beispiele dazu.%

Damit auch keine Fehler passieren, sollte man am Ende einer Zeile immer ein Prozentzeichen "\%" setzen. Die Begründung dazu findet man \href{https://tex.stackexchange.com/questions/7453/what-is-the-use-of-percent-signs-at-the-end-of-lines#7459}{hier}.%

Falls man also innerhalb eines Kapitels eine Einrückung haben möchte, muss man nur eine Leerzeile(siehe \ref{tab: leerzeile}) einfügen.\\%
\begin{table}[ht]%
\begin{tabu} to \textwidth {X[c]X[c]}%
\toprule%
	\begin{minipage}[t]{0.4\textwidth}% minipage wird genutzt, damit die Befehle ohne Einrückung geschrieben werden können
	Erste Zeile Latex Code\\%
	\\%
	Zweite Zeile Latex Code%
	\end{minipage}%
&%
	\begin{minipage}[t]{0.4\textwidth}%
	Erste Zeile Latex Code\\%
	\hspace*{10.95pt}Zweite Zeile Latex Code% indent setzen - durch \the\parindent kann man die aktuelle Größe im Dokument ins PDF schreiben lassen
	\end{minipage}\\%
\bottomrule%
\end{tabu}%
\caption{Einrückung durch Leerzeile}%
\label{tab: leerzeile}%
\end{table}\\%
%
Einen erzwungenen Umbruch sollte man nur im äußersten Notfall gebrauchen! Dieser wird wie in \ref{tab: backslash} erzwungen, erstellt eine neue Zeile \emph{(!)} und bricht den Blocksatz ab \emph{(!)}.\\%
\begin{table}[h]%
\begin{tabu} to \textwidth {X[c]X[c]}%
\toprule%
	\begin{minipage}{0.4\textwidth}%
	Erste Zeile Latex Code\textbackslash \textbackslash Zweite Zeile Latex Code
	\end{minipage}%
&%
	\begin{minipage}{0.4\textwidth}%
	Erste Zeile Latex Code\\%
	\noindent Zweite Zeile Latex Code%
	\end{minipage}\\%
\bottomrule%
\end{tabu}%
\caption{Erzwungener Zeilenumbruch}%
\label{tab: backslash}%
\end{table}\\%
Befolgt man diese Ratschläge, muss man sich so gut wie keine Gedanken machen, um die Formatierung.
%
%
\section{Dokumentstruktur}%
Oft will man Kapitel oder andere Teile der Arbeit umstrukturieren. Dafür empfiehlt es sich zumindest die Kapitel in \emph{einzelne} Dateien zu packen, die man in einer Hauptdatei zusammenfügt. Einfügen solltet ihr \textsc{NUR} nach dem Muster in Tabelle \ref{tab: einfuegen}. Der Befehl \textbackslash input\{\} wirkt dabei so, als ob ihr den Code direkt in euer Dokument einfügen würdet.%
\begin{table}[h]%

\begin{tabu} to \textwidth {X[c]X[c]}%
\toprule%
\textbackslash input\{./content/kapitelname.tex\} & Fügt den Tex-Code aus der Datei "kapitelname.tex" genau an der Stelle ein, wo der Befehl steht. Dabei liegen die Dateien im Ordner "/content/" in dem Ordner mit der Datei \emph{main-htw-abschlussarbeit.tex}\\%
\bottomrule%
\end{tabu}%
\caption{Dokumentstruktur durch Einfügen}%
\label{tab: einfuegen}%
\end{table}%
%
%
\section{Kapitelstruktur}%
Mit folgenden Befehlen aus \ref{tab: kapitelstruktur} Strukturiert man seine Überschriften. Diese 6 Ebenen solltet ihr maximal verwenden.\\%
\begin{table}[h]%
\begin{tabu} to \textwidth {X[c]X[c]}%
\toprule%
\textbackslash chapter\{Name\} & 1. Name\\%
\textbackslash section\{Name\} & 1.1. Name\\%
\textbackslash subsection\{Name\} & 1.1.1 Name \\%
\textbackslash subsubsection\{Name\} & 1.1.1.1 Name \\%
\textbackslash paragraph\{Name\} & 1.1.1.1.1 Name \\%
\textbackslash subparagraph\{Name\} & 1.1.1.1.1.1 Name \\%
\bottomrule%
\end{tabu}%
\caption{Kapitelstruktur}%
\label{tab: kapitelstruktur}%
\end{table}%
Anmerkungen zur Struktur:\\%
\begin{itemize}%
\item Außer den voreingestellten nicht nummerierten Kapiteln werden alle Kapitel nummeriert in einer Abschlussarbeit! Eine Ausnahme bildet der Anhang - weitere Informationen findet ihr im \ref{chap: anhang}%
\item Ein Unterpunkt benötigt mindestens zwei Einträge! Auf \textbackslash chapter\{\} müssen also mindestens zwei \textbackslash section\{\} folgen, damit ein Einfügen von \textbackslash section\{\} berechtigt ist!%
\item Änderungen der Kapitelstruktur sind \emph{sofort} im linken Übersichtsbereich von \emph{TexMaker} unter \textsc{structure} sichtbar, yay!%
\item Bei Änderungen der Kapitelstruktur muss man \emph{zwei Mal kompilieren}, damit man die Änderung auch im PDF sieht!%
\end{itemize}%
%
%
\section{Akzente setzen}%
Für längere Akzente benutzt man \textit{kursive Schrift}. Funktioniert wie hier in \ref{tab: kursiv} zu sehen:%
%
\begin{table}[h]%
\begin{tabu} to \textwidth {X[c]X[c]}%
\toprule%
\textbackslash textit\{kursiv\} & \textsc{kursiv} \\%
\bottomrule%
\end{tabu}%
\caption{Kursiv}%
\label{tab: kursiv}%
\end{table}%
%
Weil \textbf{fetter Text} mit der Schriftart \emph{Arial} innerhalb eines Textes wirklich fies aussieht und \underline{unterstreichen} mindestens ebenso schäbig aussieht, verwenden wir weder fetten Text noch unterstreichen wir Wörter. Stattdessen nutzen wir \textsc{Kapitälchen} und \emph{Akzente}. \textsc{Kapitälchen funktionieren} wie in \ref{tab: kapitälchen} zu sehen:\\%
\begin{table}[h]%
\begin{tabu} to \textwidth {X[c]X[c]}%
\toprule%
\textbackslash textsc\{Kapitälchen\} & \textsc{Kapitälchen} \\%
\bottomrule%
\end{tabu}%
\caption{Kapitälchen}%
\label{tab: kapitälchen}%
\end{table}%
%
Um einzelne \emph{Akzent} zu setzen, nutzt man den Befehl aus \ref{tab: akzent}. Man beobachte, dass der \emph{Akzent} in kursiver Schriftumgebung anders gesetzt wird:%
\begin{table}[h]%
\begin{tabu} to \textwidth {X[c]X[c]}%
\toprule%
\textbackslash emph\{Akzent\} & \emph{Akzent} \\%
Ein \emph{Akzent} in normaler Schriftumgebung & \textit{Ein \emph{Akzent} in kursiver Schriftumgebung}\\%
\bottomrule%
\end{tabu}%
\caption{Akzente}%
\label{tab: akzent}%
\end{table}%
%
Für Tabellenüberschriften benötigt man gelegentlich doch eine fette Schriftart. In Tabelle \ref{tab: fett} findet man den Befehl dazu. Kapitelüberschriften und einige andere Formatierungen sind übrigens ebenso von dieser Regel ausgenommen.%
\begin{table}[h]%
\begin{tabu} to \textwidth {X[c]X[c]}%
\toprule%
\textbackslash textbf\{fett\} & \textbf{fett} \\%
\bottomrule%
\end{tabu}%
\caption{Fette Schrift}%
\label{tab: fett}%
\end{table}%
%
%
\section{Besondere Zeichen}%
Quotation, da sonst Fehler entstehen, eventuell csquotes?\\%
wichtige sonderzeichen listen (zum beispiel textbackslash usw.)\\%
restliche zeichen sind in TexMaker zu finden.%%
%
%
\chapter{Bilder, Tabellen, Formeln und ihre Referenzen}%
%
%
\section{Label und Referenzen}%
\label{sec: labandref}%
%
%
\section{Tabellen}%
\label{sec: tables}%
%
%
\section{Bilder}%
Bilder müssen wie auch Tabellen zum sinnvollen Beschriften und Referenzieren in eine Umgebung geschrieben werden. Nur so kann man Bilder referenzieren (siehe \ref{sec: labandref} und sie wie auch Tabellen benennen (siehe \ref{sec: tables}.%

\begin{figure}[h]%
    \includegraphics[width=0.25\textwidth]{stura-logo-htw.pdf}%
    \caption{Logo des STURA der HTW Dresden}%
    \label{fig: stura-logo}%
\end{figure}%
dann noch Bilder:
Diese müssen in der Diplomarbeit - egal wie winzig sie sind - eine Bildunterschrift haben und ebenso in einer eigenen Zeile stehen!!! Keinen Umlaufenden Text verwenden.
%
%
\section{Positionieren von Bildern und Tabellen}%Vielleicht auch von Formeln? uff, doch noch viel zu tun.
\section{Formeln}%
Die wichtigsten Dinge für Mathematik im Dokument sind folgende:
aus\\
\begin{tabular}[h]{cc}
\hline 
 a=b  & \textbackslash begin \{equation\} a=b \textbackslash end \{equation\} \\ 
\hline 
\end{tabular} 

\begin{lstlisting}
\begin{equation}
a=b
\end{equation}
\end{lstlisting} 
wird
\begin{equation}
a=b
\end{equation}
aus 
\begin{lstlisting}
\begin{equation*}
a=b
\end{equation*}
\end{lstlisting}
wird
\begin{equation*}
a=b
\end{equation*}
ohne Nummerierung

Mathe im Text schreibt man mit dem Dollarzeichen
\begin{lstlisting}
$E = mc^2$
\end{lstlisting}
und man erhält somit $E=mc^2$.

Damit sollten die wichtigsten Arten für Mathematik in der Diplomarbeit abgedeckt sein.
%
%
\section{Literatur}%
Die Literatur schreibt man in die Datei \emph{literatur.bib}, welche relativ zu eurem Template im Verzeichnis \emph{/frontbackmatter/} liegt.%

Hier kommt euch \emph{TexMaker} zu Hilfe: Unter \emph{Bibliography -> biblatex} findet man Verschiedene Arten von Literatur, die man mit einem Klick der Datei hinzufügen kann. Dabei sollte man beachten, dass zwischen den einzelnen Quellen unbedingt eine \textsc{Leerzeile} sein sollte, da es sonst eventuell zu Problemen beim Referenzieren kommt. Anschließend kann man mit der Option \emph{Bibliography -> Clean} alle optionalen Felder automatisch entfernen lassen, sofern die Buchstaben \emph{OPT} am Anfang der Zeile nicht entfernt wurden.%

Im ersten Feld gibt man sein Label für diese Quelle an - alle anderen Felder sind selbsterklärend.%
%
%
\section{Zitieren}%
Nachdem man seine Quellen dem Literaturverzeichnis hinzugefügt hat und diese mit \emph{Labeln} versehen hat, kann man diese anschließend Zitieren.%

Wörtliche Zitate mittels \textbackslash begin  \{hyphendisplayquote \} starten, damit Silbentrennung aktiviert wird - dieses Zitat muss noch konfiguriert werden usw. siehe CS quotes.%

Damit die Literaturreferenzen im Dokument erkannt werden, muss man vor dem Kompilieren den selbst erstellten \emph{User Command} für \emph{biblatex} starten und anschließend die PDF mindestens zwei mal Kompilieren, damit das Ergebnis vollständig sichtbar wird.%
%
%
\section{Fußnoten}%
Fußnoten setzt man wie in \ref{tab: fussnoten} zu sehen.\\%
\begin{table}[h]%
\begin{tabu} to \textwidth {X[c]X[c]}%
\toprule%
Wort \textbackslash footnote\{Fußnotentext\} & Wort \footnote{Fußnotentext} \\%
\bottomrule%
\end{tabu}%
\caption{Fußnoten}%
\label{tab: fussnoten}%
\end{table}%
%
%
\section{Absatztest}

\gls{latex}
\Gls{formula}



Let's cite! Einstein's journal paper \cite{einstein} and the Dirac's 
book \cite{dirac} are physics related items. 
Next, \textit{The \LaTeX\ Companion} book \cite{latexcompanion}, the Donald 
Knuth's website \cite{knuthwebsite}, \textit{The Comprehensive Tex Archive 
Network} (CTAN) \cite{ctan} are LaTex related items; but the others Donald 
Knuth's items \cite{knuth-fa,knuth-acp} are dedicated to programming. 




Lorem ipsum dolor sit amet, consectetur adipiscing elit. Morbi convallis, massa consectetur vulputate scelerisque, est purus cursus massa, sit amet auctor nunc augue eget ex. Nunc ultricies nisl tellus, eget molestie odio vestibulum quis. Phasellus rhoncus tempus quam, imperdiet imperdiet nulla ultrices non. Nulla semper venenatis lobortis. Nam sed lectus consectetur, dignissim mi vitae, rhoncus urna. Nulla facilisi. Pellentesque purus purus, pretium quis odio vitae, mattis tristique quam. Phasellus vel tortor felis. Ut efficitur massa sit amet tellus finibus malesuada. Nam sit amet commodo eros. Phasellus tempor augue a lobortis luctus. Etiam lobortis est eget magna dignissim accumsan. Curabitur non aliquet neque. Etiam neque orci, ultricies quis ultrices vel, lobortis sed tellus. Mauris volutpat non tellus sit amet vulputate. Praesent tempus dolor vel lorem ullamcorper, id elementum libero efficitur.

Sed mi mi, eleifend sit amet tincidunt ac, eleifend quis justo. Etiam dictum iaculis justo, sit amet sagittis nunc luctus non. Etiam id aliquam ante. Maecenas gravida aliquet erat, vel suscipit urna ornare a. Etiam at ligula tortor. Praesent semper libero nisi, non ornare leo aliquet tempus. Ut maximus pharetra elementum. Integer varius hendrerit sapien, at consequat mauris mollis vel. Interdum et malesuada fames ac ante ipsum primis in faucibus. Aenean aliquam pharetra est nec iaculis. Nullam lectus quam, laoreet nec libero eget, rhoncus molestie augue.

In hac habitasse platea dictumst. Nam tempor ornare est, ac tincidunt nunc interdum et. Class aptent taciti sociosqu ad litora torquent per conubia nostra, per inceptos himenaeos. Sed volutpat magna in magna laoreet ultricies. Integer nec pellentesque nibh. Suspendisse porttitor felis id aliquet iaculis. Vivamus sodales quam nulla, quis varius dui pharetra et. Phasellus blandit ac erat ac malesuada. Etiam luctus, libero nec condimentum accumsan, nisl orci posuere quam, feugiat lacinia massa enim sed velit.\\
Cras neque libero, maximus vitae lacus non, pulvinar convallis arcu. Morbi a augue sollicitudin, ullamcorper lacus et, ultricies arcu. In nec tempus ex. Suspendisse potenti. Nullam laoreet lacus in ante malesuada, id dictum libero laoreet. Sed rhoncus volutpat massa ut accumsan. Nulla at rhoncus elit. Nam blandit consequat ipsum, et commodo erat malesuada quis. Mauris suscipit hendrerit nisl et malesuada. Phasellus ac sem sapien. Suspendisse at arcu finibus, maximus tortor a, efficitur dui. Nam volutpat purus eget dignissim consectetur. Phasellus id egestas orci. Lorem ipsum dolor sit amet, consectetur adipiscing elit. Quisque id dolor nec elit consectetur commodo non et risus. Vestibulum congue lacinia urna, quis ornare metus ultricies eu.\\
Vestibulum condimentum sagittis ipsum, a dictum justo. Integer quis magna fringilla, consectetur odio at, tristique elit. Etiam tellus purus, scelerisque nec mollis ut, feugiat ac leo. Proin libero ex, ullamcorper eget vulputate vitae, molestie a mi. In vestibulum pulvinar massa, sit amet volutpat enim feugiat vitae. Pellentesque habitant morbi tristique senectus et netus et malesuada fames ac turpis egestas. Morbi hendrerit mi hendrerit odio ultricies pulvinar. Sed eu mattis enim.
\newpage
%
%
%
% Literatur- und Quellenverzeichnis
Literatur- und Quellenverzeichnis
\newpage
%
%
% Eidesstattliche Erklärung
\addcontentsline{toc}{chapter}{Erklärung über die eigenständige Erstellung der Arbeit}%
\chapter*{Erklärung über die eigenständige Erstellung der Arbeit}%
Hiermit erkläre ich, \makeatletter \@author \makeatother , dass ich die vorgelegte Arbeit mit dem Titel \makeatletter \textsf{\@title} \makeatother selbstständig verfasst, keine anderen als die angegebenen Quellen und Hilfsmittel benutzt sowie alle wörtlich oder sinngemäß übernommenen Stellen in der Arbeit als solche und durch Angabe der Quelle gekennzeichnet habe. Dies gilt auch für Zeichnungen, Skizzen, bildliche Darstellungen sowie für Quellen aus dem Internet.\\%
Mir ist bewusst, dass die Hochschule für Technik und Wirtschaft Dresden Prüfungsarbeiten stichprobenartig mittels der Verwendung von Software zur Erkennung von Plagiaten überprüft.\\%
Ferner gestatte ich der Hochschule für Technik und Wirtschaft Dresden, die beiliegende Diplomarbeit unter Beachtung insbesondere urheber-, datenschutz- und wettbewerbsrechtlicher Vorschriften für Lehre und Forschung zu nutzen.
Es ist mir bekannt, dass für die Weitergabe oder Veröffentlichung der Arbeit die Zustimmung der HTW Dresden sowie der an der Aufgabenstellung und Durchführung der Arbeit unmittelbar beteiligten Partnereinrichtungen erforderlich ist.\\%

\vfill%

Ort, \today \hspace{35mm} \hrulefill \\
\hspace*{83mm} Unterschrift
\newpage%
%
% 
% Anlagenverzeichnis
Anlagenverzeichnis
\newpage
% Anlagen
\appendix
\newpage
% 2. Leerseite
%\afterpage{\blankpage}
\newpage
\end{document}
%Vorgaben HTW Dresden:
%Schrift:
%10pt Arial/Helvetica
%Blocksatz
%fußnoten 8pt Arial oder Helvetica (Corporate Design)
%Satzspiegel
%Ein- oder Zweiseitig
%Bindekorrektur 10mm bis 100 Seiten, 15mm bis 150 Seiten (ungefähr)
%zum außenrand: 25mm
%zur bindekorrektur: 30mm
%kopf 20mm
%fuß 20mm
%
%Absatzabstand 6pt oder Einrücken der ersten Zeile
%zeilenabstand 1-1,5 Zeilen
%
%Nummerierung erst mit Römischen Zahlen, dann arabische Zahlen
%Seitenzahl bei einseitigem Dokument (rechter Seitenspiegel) immer rechts unten
%bei zweiseitigem Dokument rechts unten bei rechtem Seitenspiegel und links unten bei linkem Seitenspiegel

%Kopf- und Fußzeile sind sichtbar zu trennen! z.B. mit Trennlinie
%farbige Abstufungen ausschließlich in grau
%Kopfzeile enthält Kapitelnummer und Kapitelname zentriert
%wird erst ab Seite 2 eines Kapitels aktiv
